% généré par scripts/06_textes.py
\dosmeta{Année de classement : 385 (directe : nom consulaire lu). Références : ICUR I 1441 = ICVaticano p 350 = ILCV 4460 = AE 1997, 166. Rattachement selon l’EDB : \emph{n.d.}. Conservation selon l’EDB : \emph{Musei Vaticani: Lapidario Cristiano}. Support : Plaque de marbre. 41 × 52 cm (h. × l.). Inscription gravée. Signes décrits par l’EDB : chrisme et chrisme avec alpha et oméga. Reproductions signalées par l’EDB : reproduction, Perret 1851 (droits réservés ; non reproduites). Fiche EDB : Filippo Piazzolla, première version 2015. Notice du catalogue p.~\pageref{not:275}.}
\txtbloc{Texte transmis (EDCS, état 2022)}{\lnum{1}Civem(!) Armeniacum Cappadocem\par
\lnum{2}n⸢o⸣mine Quirillus pius omni\par
\lnum{3}bus amicatus vixit annos\par
\lnum{4}p(lus) m(inus) LXX praecessit ad pacem\par
\lnum{5}VI Id(us) Mart(ia)\par
\lnum{6}s\par
\lnum{7}post conss(ulatum ) Rico\par
\lnum{8}mere(s) et Clearc(h)⸢i⸣ die Lunae\par
\lnum{9}bisomu(m) cum Ype« r »echiam(!) co⸢m⸣parem(!)}
\txtbloc{Texte de l’EDB (EDB31206)}{\lnum{1}((monogramma Christi cum alpha et omega))\par
\lnum{2}civem Armeniacum Cappadocem\par
\lnum{3}numine Quirillus pius omni/bus amicatus vixit annos\par
\lnum{4}ṕ(lus) ḿ(inus) LXX praecessit ad pacem\par
\lnum{5}VI  id̅(us) mart(ia)`s´ post cons̅s̅(ulatum) Rico/mere et Clearco die lunae\par
\lnum{6}bisomu cum Ype`r´echiam conparem}
